% !TeX root = ../main.tex
\section{Discussion}
\label{sec:discussion}

\subsection{Mechanism analysis of performance gains}
\label{subsec:mechanism}

The experimental results show that the proposed room-level representation outperforms the component-based baseline, particularly in precision (+17.5\%). This gain can be attributed mainly to the topological constraints inherent in the room-based graph structure. Unlike component-based models, which predict wall segments as independent entities and therefore often produce fragmented or ``floating'' walls that violate architectural logic, the room-based formulation enforces strict boundary adherence. When each room is treated as the atomic prediction unit, the model embeds a structural prior that shear walls must coincide with room boundaries. This prior filters out spurious predictions in non-boundary regions and directly reduces false positives.

Furthermore, the ablation study highlights the critical role of the FiLM conditioning mechanism. The sharp decline in density agreement ($S_{den}$) when removing conditioning suggests that geometric features alone are insufficient to drive density-aware generation. In deep GNNs, global condition signals (e.g., a simple density scalar) can weaken during recursive message passing if they are only concatenated with high-dimensional node features. In this context, conditional collapse means that the model gradually produces similar layouts across different condition inputs. The effectiveness of FiLM lies in its multiplicative modulation, which recalibrates feature activations at each layer according to the target design condition \cite{perezFiLMVisualReasoning2018}. This helps maintain sensitivity to the specified seismic intensity or height requirements throughout inference.

\subsection{Practical implications and scope}
\label{subsec:implications}

Beyond the numerical improvements reported in Section~\ref{sec:experiments}, the main practical implication of the framework lies in how it can be inserted into the preliminary design workflow. In current practice, early shear wall layout development is usually an iterative coordination process between architects and structural engineers, and the speed of that loop directly limits how many candidate schemes can be explored. The proposed method is intended for this stage rather than for final structural detailing. Based on the present dataset experiments and the two FE case studies, its current practical value should be interpreted as the ability to generate multiple boundary-feasible alternatives under different density conditions for rapid preliminary comparison. Whether such alternatives are suitable for direct project adoption still requires project-specific structural analysis and engineering review.

Regarding the scope of application, this study focuses on high-rise residential buildings in seismic regions, a typology where shear wall layouts are strictly governed by room partitions. While the underlying graph formulation is generic, the learned patterns are specific to this domain. Extending this approach to public buildings (e.g., offices or malls) with large open spaces would require adapting the room definition, potentially using virtual partitions or grid-based subdivisions, and retraining the model on domain-specific datasets. Nevertheless, the core premise of aligning structural predictions with functional spaces remains applicable.

\subsection{Limitations}
\label{subsec:limitations}

Despite the promising results, the limitations of the current framework are best understood in terms of their impact on engineering use. The first limitation lies in the geometric abstraction and applicable scope. The current room-level graph formulation is most suitable for regular high-rise residential, hotel, and apartment-type plans with predominantly rectangular or rectilinear spaces. For highly irregular layouts, large public spaces, complex corridor systems, or plans containing many fragmented spatial regions, the decomposition into rectangular room nodes may introduce virtual partitions that do not correspond to actual walls or structural design decisions. Although the constraint mask prevents shear wall placement on these artificial boundaries, the additional nodes and edges may still influence message passing and reduce the effectiveness of the learned representation. In such scenarios, component-level graph representations may be advantageous because they preserve local segment-level topology, small openings, and detailed geometric alignment more directly.

The second limitation concerns optimization target mismatch. The model learns from existing engineering layouts and is regularized by density-related constraints, but it is not trained to optimize structural response indicators such as drift ratio, period, or torsional behavior. The FE case studies show that generated schemes can remain close to engineer-designed solutions under the tested settings, yet this should be interpreted as feasibility evidence rather than proof of structural optimality. Integrating differentiable or surrogate structural analysis into training would be a more direct route toward performance-oriented generation.

The third limitation is the granularity of conditional control. The current model uses three discrete condition groups derived from seismic intensity and building height. This setting is consistent with the available data but does not provide continuous control over engineering parameters such as PGA, building height, structural period, drift demand, or torsional behavior. Therefore, the present framework should be interpreted as discrete density-controlled preliminary generation rather than fully continuous performance-conditioned structural design. Expanding the dataset and moving from discrete categories to continuous condition descriptions would be necessary for broader deployment.
